\section{AR4 Restoration and Hardware Modifications}
The AR4 robotic manipulator was handed over to the current
team by the previous graduation project team after their completion of the project.
At the time of handover, the robot was delivered in a fully assembled
condition and was generally operational. The unit is an AR4 MK3 model designed by
Chris Annin. However, due to the duration and intensity of prior
usage during development and testing, several components exhibited
wear and assembly-related issues. As a result, multiple mechanical
and electrical corrective actions were required before the 
robot could be considered reliable for accurate motion execution and 
experimental work.

\subsection{Joint 3 Mechanical Transmission Upgrade}
Joint~3 of the AR4 robot exhibited significant backlash due 
to wear and limitations in the original mechanical transmission 
design. To address this issue, the joint transmission system was 
completely replaced in accordance with the official upgrade 
provided by Chris Annin. The upgrade consisted of a redesigned 
gear set and a new timing belt, both intended to improve torque 
transmission and reduce mechanical play within the joint.

\begin{figure}[H]
    \centering
    \includegraphics[width=0.75\textwidth]{figures/chapter3/AR4 modifications/Joint3_mods.png}
    \caption{Upgraded mechanical transmission of Joint~3 }
    \label{fig:joint3_upgrade}
\end{figure}

After implementing the upgraded components, the backlash in 
Joint~3 was reduced from approximately $6^\circ$ to $3^\circ$. 
This reduction resulted in smoother joint motion and improved 
positioning accuracy, directly enhancing the overall kinematic 
performance of the robot.

\subsection{Shaft Key Re-Manufacturing in Joint 3}
Despite the improvements achieved through the mechanical 
transmission upgrade, residual backlash was still observed in 
Joint~3. Further investigation revealed that the shaft key 
responsible for transmitting motion from the second gear stage 
to the robotic arm link exhibited wear and improper fit. 
This condition introduced mechanical play within the joint, 
negatively affecting positioning accuracy.

\begin{figure}[H]
    \centering
    
\includegraphics[width=0.45\textwidth, angle=90]{figures/chapter3/AR4 modifications/Shaft_key.png}
    \caption{Re-manufactured shaft key used in Joint~3 to eliminate mechanical backlash}
    \label{fig:shaft_key}
\end{figure}

To resolve this issue, the shaft key was re-manufactured with 
precise dimensions to ensure a tight and accurate fit between 
the gear and the arm shaft using hot-wire machining. The newly manufactured key eliminated 
relative motion between the connected components, effectively 
removing the remaining backlash in Joint~3. As a result, 
Joint~3 achieved backlash-free motion.

\subsection{Encoder Signal Noise in 5th and 6th joint}

During system testing and motion execution, unstable and noisy encoder 
readings were observed in Joints~5 and~6 of the AR4 robot. These fluctuations 
resulted in inaccurate joint position feedback, which significantly affected 
the forward and inverse kinematics calculations. Consequently, commanded joint 
positions did not always correspond to the expected physical motion of the robot.

A systematic debugging process was conducted, including signal tracing and 
verification of the encoder wiring and power connections. This investigation 
revealed that the encoders of Joints~5 and~6 were not properly connected to 
the regulated $5\,\text{V}$ power supply. The improper power connection caused 
voltage instability, leading to fluctuating encoder signals and unreliable 
position measurements.

After correcting the power supply connections and ensuring proper grounding, 
the encoder signals became stable and consistent. This correction restored 
accurate joint position feedback.

\subsection{Joint 1 Backlash Due to Missing Fasteners}

During mechanical inspection and testing of the AR4 robot,
abnormal backlash was observed in Joint~1. This backlash was estimated 
to be approximately $2^\circ$ and was found to adversely affect the base
positioning accuracy of the robot, particularly during large workspace
movements.

A detailed mechanical inspection revealed that multiple fastening screws 
within the Joint~1 assembly were missing. The absence of these fasteners 
reduced the structural rigidity of the joint, allowing relative motion 
between assembled components and introducing mechanical play during operation.

After identifying the missing screws, all required fasteners were reinstalled 
and properly tightened according to the recommended assembly specifications. 
This action restored the mechanical integrity of Joint~1 and eliminated the 
observed backlash. As a result, the positional accuracy and stability of the 
robot base joint were significantly improved.

\subsection{Summary}
Following the completion of all mechanical restorations and electrical corrections, the AR4 robot reached a stable and reliable hardware state suitable for precise operation. The corrective actions performed eliminated sources of mechanical backlash, improved structural rigidity, and ensured accurate encoder feedback across all joints. With these issues resolved, the robot was prepared for accurate kinematic modeling and motion execution.

Based on this finalized hardware condition, a comprehensive calibration procedure was subsequently carried out to align the mathematical kinematic model with the physical robot. The calibration methodology, implementation, and results are presented in the following section.
